\documentclass[11pt]{article}
\usepackage[margin=1in]{geometry}
\usepackage{times}
\usepackage{hyperref}
\hypersetup{colorlinks=true, linkcolor=black, urlcolor=blue, citecolor=black}
\title{Rotating the Electrodes Does Not Rotate the Bond Energy: Philipp Kanarev's Low-Energy Water Electrolysis}
\author{Flyxion \\ \small Independent Researcher}
\date{}
\begin{document}
\maketitle

\begin{abstract}
Philipp Kanarev, a Russian engineering academic, claimed from the 1990s that a specifically configured rotating or pulsed electrolysis cell could split water into hydrogen and oxygen using substantially less energy than the water molecule's known bond dissociation energy requires, framing the result within his own broader, non-standard theoretical framework describing electron and atomic structure. This paper argues that the claimed energy input reductions trace to the same category of measurement error common across the "over-unity electrolysis" genre addressed elsewhere in this series in the Stanley Meyer entry, specifically comparing average or root-mean-square voltage and current measurements inappropriate for the highly non-sinusoidal, pulsed waveforms these cells use, and that Kanarev's broader theoretical framework, like other total replacements for standard atomic and quantum theory addressed elsewhere in this series, has not been shown to reproduce the extensively confirmed predictions of standard quantum chemistry regarding molecular bond energies that any competing framework must match before its novel predictions can be taken seriously.
\end{abstract}

\section{Introduction}

Philipp Kanarev, a professor of mechanics at a Russian technical university, published extensively from the 1990s onward claiming that specially designed electrolysis cells, using pulsed direct current and, in some variants, rotating electrode assemblies, could dissociate water into hydrogen and oxygen gas using energy input substantially below the approximately 286 kilojoules per mole that standard thermochemistry establishes as the water molecule's dissociation enthalpy, attributing the result to his own independently developed theoretical model of atomic and subatomic structure, which departs significantly from standard quantum mechanics in ways Kanarev argued permitted this apparent efficiency gain.

\section{The Measurement Problem With Pulsed Electrolysis Power}

Determining the true electrical power delivered to an electrolysis cell driven by a non-sinusoidal, pulsed voltage and current waveform requires either direct calculation of the true root-mean-square power, integrating the instantaneous product of voltage and current across the actual waveform shape, or a power meter specifically designed to handle non-sinusoidal waveforms correctly; using a simple average-reading voltmeter and ammeter, or applying formulas valid only for steady direct current or sinusoidal alternating current to a sharply pulsed waveform, can substantially underestimate the true delivered electrical power, in some documented cases in the broader pulsed-electrolysis literature by a factor of two or more, an error category independently identified and specifically addressed in critical technical reviews of Kanarev's published measurement methodology by electrical engineers examining his specific reported power figures and finding the same waveform-measurement mismatch present.

\section{Why This Fully Accounts for the Claimed Anomaly Without Requiring New Physics}

Once the true, correctly integrated electrical power delivered to Kanarev's pulsed cells is calculated from the actual waveform rather than from an average-reading meter inappropriate to it, independent reanalyses of his published data, using his own reported waveform characteristics, have found energy input consistent with, or exceeding, the standard water dissociation energy requirement, fully accounting for the observed hydrogen and oxygen output without requiring any departure from established thermochemistry, a finding that mirrors the resolution of the Stanley Meyer water fuel cell claim addressed elsewhere in this series, where a comparable pulsed-waveform measurement issue similarly explained an apparent efficiency anomaly.

\section{The Theoretical Framework Faces the Same Bar as Other Replacement Theories in This Series}

Kanarev's broader theoretical program, proposing a substantially revised model of atomic structure and electron behavior distinct from standard quantum mechanics, is, like the space vortex theory addressed in the Tewari entry and the classical hydrogen model addressed in the hydrino entry elsewhere in this series, a candidate replacement for an extraordinarily well confirmed existing theoretical framework, standard quantum chemistry's account of molecular bond energies, which has been validated to high precision across an enormous range of spectroscopic, thermochemical, and computational chemistry results. Any replacement theory must, at minimum, be independently shown to reproduce this existing body of confirmed prediction before its novel claims in an untested regime, reduced water dissociation energy under specific electrical conditions, can be evaluated as more than speculative extrapolation from an unvalidated framework, a demonstration Kanarev's published work has not been shown, by independent quantum chemists reviewing it, to have satisfied.

\section{The Entropic Gravity Framing}

Kanarev's claim concerns chemical bond dissociation energy rather than gravitational coupling directly, and is included in this series as a comparative case in the pulsed-electrolysis over-unity genre alongside the Stanley Meyer entry, illustrating that the same waveform power-measurement error recurs independently across different inventors and different specific theoretical justifications offered for the resulting apparent anomaly.

\section{Objections}

Supporters argue that Kanarev's status as an academic engineering professor with an extensive published body of work, distinct from the less formally credentialed inventors surveyed elsewhere in this series, warrants a different evidentiary standard or greater initial credibility. An academic credential and an extensive publication record, much of it in journals and venues outside mainstream peer-reviewed physics and chemistry, do not substitute for independent replication and correct measurement methodology, and the specific waveform power-measurement critique raised above addresses Kanarev's own published data and methodology directly rather than relying on his credentials or venue of publication as grounds for skepticism.

\section{Conclusion}

Kanarev's claimed reduced-energy water electrolysis traces to a well documented pulsed-waveform power measurement error also identified in the Stanley Meyer water fuel cell claim, correctly calculated power delivery is consistent with standard water dissociation thermochemistry, and his broader replacement theory of atomic structure has not been independently shown to reproduce the extensively confirmed predictions of standard quantum chemistry that any competing framework must match before its own novel claims can be taken as more than speculative.

\end{document}
